\section{Methods}
\label{sec:methods}

\subsection{Study population and prediction time}
We used electronic health records from six Israeli general hospitals in the Kineret network for this retrospective study. The hospitals share a data warehouse based on the Observational Medical Outcomes Partnership common data model. Before abstraction, local measurement, medication, and observation codes were mapped to a shared clinical vocabulary and measurement units were harmonized \citep{oded2027kineret-jbi}.

The source cohort comprised adults admitted for 48 hours to 14 days with at least one overnight measurement (18:00--07:00). An admission qualified if the patient had a documented diabetes diagnosis or met either glucose criterion in the first 48 hours. The glucose criteria were two readings of at least 180 mg/dL taken 2--24 hours apart, or one reading of at least 250 mg/dL. The source study excluded admissions to pediatric and obstetric departments, deaths within 48 hours, and admissions with a duplicated visit identifier assigned to more than one person \citep{oded2027kineret-jbi}.

We chose 96 hours as the primary prediction time to allow repeated measurements and responses to treatment to enter the history. Both the raw and abstracted records had to extend beyond that point. The calendar dates came from the approved multicenter protocol for evaluating the national inpatient diabetes recommendations, issued in September 2023. That study drew on a period before implementation (August 2022--July 2023) and a later period allowing time for adoption (August 2024--December 2025) \citep{hochberg2023inpatient,oded2027kineret-jbi}. The clinical study evaluated the change in adherence across these periods. For prediction, the same calendar eligibility rule applied to training, validation, and testing. We required admission on or after 1 August 2022 and a full 14-day horizon by 31 December 2025. We kept the original split by person (70\% training, 15\% validation, 15\% testing), so repeat admissions stayed together. Within these dates, 20,994 training, 4,427 validation, and 4,508 test admissions reached the four-day cutoff.

\subsection{Outcomes and timing}
The forecast period began immediately after hour 96 and ended at hour 336 (day 14), inclusive. An admission could count as positive for several outcomes. For every complication, including infection, kidney and cardiovascular complications, an earlier occurrence did not remove the admission from the risk set. A positive label required an event timestamp in the forecast period. For example, hypoglycemia both before and after hour 96 counted as positive. An episode starting before the cutoff counted again only if the event table contained a subsequent occurrence. For each outcome that occurred, the timing target was the hour of its first occurrence within the forecast period, measured from admission. Total length of stay ran from admission to the recorded visit end, using the release marker when the visit-end field was missing.

Each admission had one terminal marker at the recorded visit end, either \textsc{release} or \textsc{death}. A linked death within 30 days after discharge was encoded as \textsc{death} at visit end. The linkage also accepted death records dated up to 14 days before administrative discharge to capture in-hospital deaths. For mortality, the forecast therefore concerned whether the death marker fell within the prediction window. For example, discharge on day 8 followed by a linked death on day 25 counted as a positive mortality label at day 8. Thus the 14-day horizon applied to the visit end, while ascertainment of death extended beyond discharge.

Of ten candidate outcomes, eight occurred in at least 1\% of admissions in the original day-four training set. We retained hyperglycemia, hypoglycemia, their severe forms, kidney complications, cardiovascular complications, infection, and death. The outcomes remained fixed after date filtering, when the training prevalence of infection fell to 0.991\%.

Hyperglycemia required a sustained elevation at or above 180 mg/dL or a single reading at or above 250 mg/dL. The latter also qualified as severe hyperglycemia. Hypoglycemia required sustained glucose at or below 70 mg/dL or a single reading at or below 54 mg/dL, the severe threshold. The kidney endpoint drew on creatinine measurements and recorded kidney events. The infection definition combined a change in temperature or white cell counts with diagnostic testing and antibiotic treatment. We used the outcome criteria described for the Kineret cohort \citep{oded2027kineret-jbi}. We reconciled event times across the raw and abstracted records to give every model the same outcome labels.

\subsection{Clinical histories and model comparisons}
\label{sec:methods-ladder}
Seven model configurations used the same admissions, outcomes, and prediction times. We first compared INTERVenE-Enc with models that used the original observations. We then kept the encoder architecture fixed and varied how the intervals were defined.

\emph{Models of individual observations.} Logistic regression used summaries of each measured variable, including its first and last values, distribution, slope, and measurement frequency. Separate linear regressions on the same features estimated onset and length of stay. The supervised STraTS model (denoted ss-STraTS) represented each observation by the variable measured, its time, and its value. Attention combined the observations into a history for each admission \citep{tipirneni2022strats}. The observation stream included vital signs, laboratory measurements such as glucose and creatinine, and recorded treatments including insulin. Both models also received admission-level covariates. Supplementary Methods A gives the construction of the regression vector.

\emph{Interval models.} In the clinical knowledge-base configuration (KB), the Mediator converted the record into intervals denoting states, trends, events, and contexts. INTERVenE-Enc represented the start and end of each interval separately. Each boundary encoded the underlying variable, the abstraction and its value, and the elapsed time since admission \citep{oded2026intervene,kazemi2019time2vec}. Supplementary Methods B illustrates how measurements become clinical intervals and how the encoder represents their boundaries. A bidirectional encoder with four layers, two attention heads per layer, and an embedding dimension of 128 used the observed history to estimate risk, onset, and length of stay.

For the statistical intervals, we used a variant of Symbolic Aggregate approXimation (SAX) \citep{lin2003symbolic}. We averaged each numeric variable over four-hour periods, standardized these averages using the variable's training-population mean and standard deviation, and assigned each to one of seven bands defined by Gaussian quantiles. We used only periods completed by prediction time, left empty periods unfilled, and merged adjacent periods with the same symbol. The SAX encoder received these bands and raw categories. The KB encoder also had access to contextual interpretations and derived complication events. Their comparison thus measures the combined contribution of the clinical definitions and the events they make available.

\emph{Guideline assessment.} We added QA to logistic regression, ss-STraTS, and the clinical interval encoder in turn. For each of 13 named patterns, the model received both an indicator of whether the pattern occurred during the observation window and the mean of its recorded compliance scores. We included only assessments that began after or at admission and whose evaluation period had ended by the prediction cutoff. Each temporal abstraction knowledge (TAK) rule defined how long to wait for the expected action. For example, the rule for continuing basal insulin on admission allowed up to 48 hours, the upper limit of its timing trapezoid. Only once the applicable window closed could the assessment record that an expected action was absent. A pattern whose window was still open at hour 96 contributed neither a score nor an occurrence at that cutoff. The same completion rule applied to its dated pattern tokens. If a pattern never occurred, both values were zero. If it occurred and scored zero throughout, the occurrence indicator was one and the mean score was zero. The KB+QA encoder also received the dated pattern intervals, allowing it to relate a particular assessment to the surrounding clinical history. We evaluated these two outputs of the Mediator together. For meal-related patterns, the rules used assumed daily meal times because the record contained neither actual meal times nor intake \citep{oded2027kineret-jbi}.

\subsection{Model fitting and evaluation}
\label{sec:methods-eval}
We learned model weights from training patients. Validation patients were reserved for model selection and, where applicable, calibration. Test patients were reserved for evaluation. INTERVenE learned representations of interval boundaries, learned to recover masked tokens in admission histories, and then trained on the clinical outcomes. For supervised fitting, we used binary cross-entropy across the eight risk targets and added 0.1 times the mean squared timing error, standardized separately for each target. An onset contributed only when the complication occurred, and we averaged timing losses across targets with observed times. We used AdamW with a head learning rate of $10^{-4}$ and selected the checkpoint by validation loss. During supervised training, an admission contributed one example at each daily cutoff, from day one to day thirteen, for which its record continued beyond the cutoff. For example, a record ending at 120 hours contributed four training examples, using the first 24, 48, 72, and 96 hours respectively. We restricted the history at each cutoff, excluded terminal discharge and death markers, and determined which outcomes occurred between the cutoff and day 14. Admission-level inputs included elapsed observation time and time remaining to day 14. At evaluation, each eligible test admission contributed one set of day-four predictions to the primary comparison.

Our primary measure was the support-weighted area under the precision--recall curve (AUPRC). We calculated average precision for each outcome, then weighted it by the number of admissions positive for that outcome. We also reported macro AUPRC, which weights the outcomes equally, and area under the receiver operating characteristic curve (AUROC). For timing, we calculated mean absolute error (MAE) in hours, using cases of each complication for onset and admissions with a recorded duration for length of stay. We assessed probability estimates using reliability plots, the Brier score, and expected calibration error across 15 bins. Where validation predictions were available, we fitted temperature scaling on those predictions and compared calibration on the test set before and after scaling.

We obtained percentile 95\% confidence intervals by resampling people 2,000 times, keeping each person's admissions together and using the same resample for both models in a comparison. For paired differences, we calculated two-sided bootstrap tests by centering the distribution of differences under the null and applying a finite-sample correction. We used the Benjamini--Hochberg procedure across outcomes within each comparison and metric, and across model comparisons within each aggregate metric and averaging scheme. An aggregate discrimination estimate required a defined value for every outcome. We flagged intervals with fewer than 95\% valid resamples and outcome-specific estimates with fewer than five positive or five negative admissions.

\subsection{Observation windows, hospitals, and recording frequency}
We evaluated the fitted sequence models at each daily cutoff from day two to day six. We compared models within each cutoff because a later prediction includes only longer admissions and leaves less time for an event to occur. Logistic regression was evaluated at the primary day-four cutoff only.

For the hospital analysis, we fitted the KB encoder six times, reserving a different hospital for testing each time. A patient seen at the test hospital contributed no admissions to training or validation, including stays at other hospitals. Supplementary Table~S7 gives performance by age, sex, and hospital for the single KB model evaluated on the pooled test set.

To characterize differences in recording practice between Kineret and the prior ICU study, we counted glucose, serum creatinine, heart rate, and temperature measurements in the first 96 hours \citep{oded2026intervene}. We used the raw inputs before abstraction in both datasets. Kineret contributed the primary test admissions. MIMIC-IV contributed all processed admissions with an explicit admission marker and a raw record extending beyond 96 hours. We counted each variable once per timestamp and divided each admission's count by four, including admissions with no measurements. Among admissions with at least two measurements of a variable, we first calculated each admission's median time between readings, then summarized those medians across admissions.

\subsection{Interpreting predictions and checking the clinical intervals}
We examined whether the trained encoder preserved relationships in the clinical vocabulary. For glucose, for example, the raw concept is the measurement, the abstractions include glucose state and glucose trend, and each takes a value such as high or decreasing. An interval with a given value has a START and an END token. With time held fixed and dropout disabled, we compared cosine distances for four kinds of token pairs. These were the two boundaries of the same valued concept, tokens sharing an abstraction, tokens sharing only the raw concept, and tokens from different raw concepts. We used uniform manifold approximation and projection (UMAP) to display these relationships and calculated all distances in the original embedding space.

We fitted a sparse transcoder to each of the four layers of the day-four KB+QA encoder \citep{dunefsky2024transcoders,rajamanoharan2024jumping}. A transcoder approximates a layer's output with learned features, few of which are active at any one position. We used 300 training patients to fit the transcoders (20 epochs, 16-fold feature expansion, up to 256 token positions per admission, seed 2023). For 64 validation patients, we assessed how much of the layer-output variance the transcoders explained ($R^2$) and what fraction of the features were active, using 13,013 positions per layer.

We then examined one admission from each of 300 test patients. For each outcome and input position, we multiplied each feature's activation by the gradient of the predicted log odds with respect to that feature and summed across features and layers. The input at that position supplied the clinical name, such as a decreasing glucose trend. We combined start and end boundaries and repeat occurrences of the same concept within an admission, then averaged signed scores across admissions, counting absent concepts as zero. We ranked concepts by the absolute value of this mean and retained the sign to indicate the direction of the contribution. Bootstrap intervals for the highest-ranked concepts were exploratory because selection and estimation used the same sample. Throughout, the encoder used its original layer outputs to make predictions. We compared its predictions with and without the attribution procedure in all 300 admissions as a separate check.

To assess the stability of the intervals, we sampled 500 people by hospital and recording density and selected one admission per person. We reran the Mediator on the first 96 hours after adding uniform noise scaled to 1\% or 5\% of local threshold spacing or removing 5\% or 10\% of observations. For 26 directly defined numeric states and their associated trends, we measured agreement as the time with matching labels divided by the time covered in either record. We also measured coverage lost or gained. We averaged repeated perturbations within person and obtained confidence intervals by weighted, stratified resampling of people.
